%%%%%%%%%%%%%%%%
%%% exod 17 %%%
%%%%%%%%%%%%%%%%

\Chapter{17}

\Verse{1}
{
	\hJ{וַיִּסְעוּ כּל עֲדַת בְּנֵי יִשְׂרָאֵל מִמִּדְבַּר סִין לְמַסְעֵיהֶם עַל פִּי יְהֹוָה וַיַּחֲנוּ בִּרְפִידִים וְאֵין מַיִם לִשְׁתֹּת הָעָם׃}
}
{
	\eJ{
		17 And all of the congregation of the children of Israel\\
		traveled from the wilderness of Sîn on their travels by\\
		YHWH’s word, and they camped in Rephidim. And there was no\\
		water for the people to drink.\\
	}
}
{
}

\Verse{2}
{
	\hJ{וַיָּרֶב הָעָם עִם מֹשֶׁה וַיֹּאמְרוּ תְּנוּ לָנוּ מַיִם וְנִשְׁתֶּה וַיֹּאמֶר לָהֶם מֹשֶׁה מַה תְּרִיבוּן עִמָּדִי מַה תְּנַסּוּן אֶת יְהֹוָה׃}
}
{
	\eJ{
		And the people quarreled with Moses. And they said, “Give\\
		us water, and let us drink.” And Moses said to them, “Why\\
		do you quarrel with me? Why do you test YHWH?!”\\
	}
}
{
%	\footnote{
%		test. In each of the two chapters that precede this, God is described as
%		testing the people (15:25; 16:4). (It is the same word that is used for
%		God’s testing Abraham in ordering him to sacrifice Isaac—Gen Footnote
%		22:1.) Now the people dare to test God, and Moses asks them why they do
%		that. It is a rhetorical question: the point is that there is no reason
%		to test God. It is presumptuous, because God can be counted upon. This
%		is significant enough that at the end of the episode the place is named
%		Massah, meaning “Test,” because of what the people dared to do there.
%	}
}

\Verse{3}
{
	\hJ{וַיִּצְמָא שָׁם הָעָם לַמַּיִם וַיָּלֶן הָעָם עַל מֹשֶׁה וַיֹּאמֶר לָמָּה זֶּה הֶעֱלִיתָנוּ מִמִּצְרַיִם לְהָמִית אֹתִי וְאֶת בָּנַי וְאֶת מִקְנַי בַּצָּמָא׃}
}
{
	\eJ{
		And the people thirsted for water there, and the people\\
		complained at Moses and said, “Why is this that you\\
		brought us up from Egypt: to kill me and my children and\\
		my cattle with thirst?!”\\
	}
}
{
}

\Verse{4}
{
	\hJ{וַיִּצְעַק מֹשֶׁה אֶל יְהֹוָה לֵאמֹר מָה אֶעֱשֶׂה לָעָם הַזֶּה עוֹד מְעַט וּסְקָלֻנִי׃}
}
{
	\eJ{
		And Moses cried to YHWH, saying, “What shall I do to this\\
		people? A little more and they’ll stone me!”\\
	}
}
{
}

\Verse{5}
{
	\hJ{וַיֹּאמֶר יְהֹוָה אֶל מֹשֶׁה עֲבֹר לִפְנֵי הָעָם וְקַח אִתְּךָ מִזִּקְנֵי יִשְׂרָאֵל וּמַטְּךָ אֲשֶׁר הִכִּיתָ בּוֹ אֶת הַיְאֹר קַח בְּיָדְךָ וְהָלָכְתָּ׃}
}
{
	\eJ{
		And YHWH said to Moses, “Pass in front of the people and\\
		take some of Israel’s elders with you, and take your staff\\
		with which you struck the Nile in your hand, and you’ll\\
		go.\\
	}
}
{
}

\Verse{6}
{
	\hJ{הִנְנִי עֹמֵד לְפָנֶיךָ שָּׁם עַל הַצּוּר בְּחֹרֵב וְהִכִּיתָ בַצּוּר וְיָצְאוּ מִמֶּנּוּ מַיִם וְשָׁתָה הָעָם וַיַּעַשׂ כֵּן מֹשֶׁה לְעֵינֵי זִקְנֵי יִשְׂרָאֵל׃}
}
{
	\eJ{
		Here, I’ll be standing in front of you there on a rock at\\
		Horeb. And you’ll strike the rock, and water will come out\\
		of it, and the people will drink.” And Moses did so before\\
		the eyes of Israel’s elders.\\
	}
}
{
%	\footnote{
%		on a rock at Horeb. The word “rock” refers not to a small stone but to a
%		crag of the mountain. God stands at a crag of Mount Horeb/Sinai.
%		Footnote 17:6. at Horeb. Here is the story’s geography: Moses and the
%		elders are to go on ahead of the people to Mount Horeb (Sinai). The
%		miracle of the water will occur there, not at the site of the quarrel
%		and testing. Only the elders see it, not the people (17:6). The people
%		receive the water somewhere near Horeb, and then they return to the camp
%		in the region of Rephidim (or perhaps the water flows all the way to
%		where the people are). The site of their quarrel there is named Massah
%		and Meribah. Amalek attacks and fights them there (17:8). Then Jethro
%		goes to meet Moses at Horeb, where Moses is now encamped (18:5). Moses
%		comes out to meet him and Moses’ wife and children there (18:7). Then
%		Aaron and the elders join them for a sacred meal there. The people come
%		there later from Rephidim (19:2). The important thing, often overlooked,
%		is that the miracle of the water from the rock takes place at Mount
%		Sinai. I learned this from William Propp. Footnote 17:6. you’ll strike
%		the rock, and water will come out of it. Moses, who was drawn from
%		water, and whose name stands for drawing from water, now draws water
%		from the rock. He also met his wife by defending her right to draw
%		water; the first plague he introduced involved turning the waters (from
%		which he had once been drawn) to blood; he initiated the splitting of
%		the sea; and he acted to sweeten the bitter water at Marah. There will
%		be several more incidents involving water in Moses’ life. It seems to be
%		a key element in his destiny, which will culminate in his tragic act
%		when he strikes a rock to get water again at a second Meribah (Numbers
%		20). For this act he will be forbidden to enter the promised land. His
%		life will end, as it began, at a river, the Jordan, which he will not be
%		permitted to cross.
%	}
}

\Verse{7}
{
	\hJ{וַיִּקְרָא שֵׁם הַמָּקוֹם מַסָּה וּמְרִיבָה עַל רִיב בְּנֵי יִשְׂרָאֵל וְעַל נַסֹּתָם אֶת יְהֹוָה לֵאמֹר הֲיֵשׁ יְהֹוָה בְּקִרְבֵּנוּ אִם אָיִן׃}
}
{
	\eJ{
		And he called the place’s name Massah and Meribah because\\
		of the quarrel of the children of Israel and because of\\
		their testing YHWH, saying, “Is YHWH among us or not?”\\
	}
}
{
%	\footnote{
%		he called the place’s name Massah and Meribah. That is, he named it
%		Testing and Quarrel.
%	}
}

\Verse{8}
{
	\hE{וַיָּבֹא עֲמָלֵק וַיִּלָּחֶם עִם יִשְׂרָאֵל בִּרְפִידִם׃}
}
{
	\eE{And Amalek came and fought with Israel in Rephidim.}
}
{
}

\Verse{9}
{
	\hE{וַיֹּאמֶר מֹשֶׁה אֶל יְהוֹשֻׁעַ בְּחַר לָנוּ אֲנָשִׁים וְצֵא הִלָּחֵם בַּעֲמָלֵק מָחָר אָנֹכִי נִצָּב עַל רֹאשׁ הַגִּבְעָה וּמַטֵּה הָאֱלֹהִים בְּיָדִי׃}
}
{
	\eE{
		And Moses said to Joshua, “Choose men for us and go, fight\\
		against Amalek. Tomorrow I’ll be standing up on the\\
		hilltop, and the staff of God in my hand.”\\
	}
}
{
%	\footnote{
%		fight against Amalek. Israel was not sent by way of the Philistines’
%		land because they would be afraid of war (13:17), but now, just a few
%		chapters later, they are sent to war against the Amalek! And Israel
%		prevails! How do we explain this? Apparently the intervening events have
%		made a change in the Israelites so that now they have the confidence and
%		fortitude to fight. What are the intervening events: Red Sea, Marah,
%		manna, Meribah. They have received four miraculous signs that their God
%		is still with them even after they have left Egypt.
%	}
}

\Verse{10}
{
	\hE{וַיַּעַשׂ יְהוֹשֻׁעַ כַּאֲשֶׁר אָמַר לוֹ מֹשֶׁה לְהִלָּחֵם בַּעֲמָלֵק וּמֹשֶׁה אַהֲרֹן וְחוּר עָלוּ רֹאשׁ הַגִּבְעָה׃}
}
{
	\eE{
		And Joshua did as Moses said to him, to fight against\\
		Amalek; and Moses, Aaron, and Hur went up to the hilltop.\\
	}
}
{
}

\Verse{11}
{
	\hE{וְהָיָה כַּאֲשֶׁר יָרִים מֹשֶׁה יָדוֹ וְגָבַר יִשְׂרָאֵל וְכַאֲשֶׁר יָנִיחַ יָדוֹ וְגָבַר עֲמָלֵק׃}
}
{
	\eE{
		And it was: when Moses would lift his hand, then Israel\\
		would predominate; and when he would rest his hand, then\\
		Amalek would predominate.\\
	}
}
{
}

\Verse{12}
{
	\hE{וִידֵי מֹשֶׁה כְּבֵדִים וַיִּקְחוּ אֶבֶן וַיָּשִׂימוּ תַחְתָּיו וַיֵּשֶׁב עָלֶיהָ וְאַהֲרֹן וְחוּר תָּמְכוּ בְיָדָיו מִזֶּה אֶחָד וּמִזֶּה אֶחָד וַיְהִי יָדָיו אֱמוּנָה עַד בֹּא הַשָּׁמֶשׁ׃}
}
{
	\eE{
		And Moses’ hands were heavy. And they took a stone and set\\
		it under him, and he sat on it, and Aaron and Hur held up\\
		his hands, one from this side and one from this side, and\\
		his hands were supported until the sunset.\\
	}
}
{
}

\Verse{13}
{
	\hE{וַיַּחֲלֹשׁ יְהוֹשֻׁעַ אֶת עֲמָלֵק וְאֶת עַמּוֹ לְפִי חָרֶב׃}
}
{
	\eE{And Joshua defeated Amalek and his people by the sword.}
}
{
}

\Verse{14}
{
	\hE{וַיֹּאמֶר יְהֹוָה אֶל מֹשֶׁה כְּתֹב זֹאת זִכָּרוֹן בַּסֵּפֶר וְשִׂים בְּאזְנֵי יְהוֹשֻׁעַ כִּי מָחֹה אֶמְחֶה אֶת זֵכֶר עֲמָלֵק מִתַּחַת הַשָּׁמָיִם׃}
}
{
	\eE{
		And YHWH said to Moses, “Write this—a memorial—in a scroll\\
		and set it in Joshua’s ears, because I shall wipe out the\\
		memory of Amalek from under the skies!”\\
	}
}
{
}

\Verse{15}
{
	\hE{וַיִּבֶן מֹשֶׁה מִזְבֵּחַ וַיִּקְרָא שְׁמוֹ יְהֹוָה נִסִּי׃}
}
{
	\eE{
		And Moses built an altar and called its name “YHWH is my\\
		standard.”\\
	}
}
{
}

\Verse{16}
{
	\hE{וַיֹּאמֶר כִּי יָד עַל כֵּס יָהּ מִלְחָמָה לַיהֹוָה בַּעֲמָלֵק מִדֹּר דֹּר׃}
}
{
	\eE{
		And he said, “Because—hand on YH’s throne—YHWH has war\\
		against Amalek from generation to generation.”\\
	}
}
{
%	\footnote{
%		hand on YH’s throne. No one is sure what this means. It may be an oath
%		formula. The word for “throne,” Hebrew ks, may be a form of kiss’, as
%		some versions have understood it.
%	}
}
